% ============================================================
%  PART V --- Settings
% ============================================================
\part{Settings}

% ── Chapter 20: Settings Overview ───────────────────────────
\chapter{Ajustes --- Sinopsis}
\label{ch:settings}

Open Settings by tapping \iconlabel{settings-iec}{Settings} in the Bottom Bar.
La configuración de escritorio e iOS se organiza en diez pestañas. Android añade un
eleventh \textbf{Android} tab after System.
Los cambios tienen efecto inmediatamente a menos que se observe explícitamente un reinicio.

\screenshot{settings-overview}{Panel de configuración con las pestañas apropiadas para plataformas visibles en la barra de pestañas}{}

\rowcolors{2}{LtGray}{white}
\begin{longtable}{C{0.6cm}L{3.0cm}L{6.8cm}L{3.3cm}}
\toprule
\rowcolor{Navy}
\color{white}\sffamily\bfseries \# &
\color{white}\sffamily\bfseries Tab name &
\color{white}\sffamily\bfseries Primary function &
\color{white}\sffamily\bfseries Restart needed? \\
\midrule
0 & Main           & Window mode, language, UI scale, comfort preset, launcher grid, coordinate format & Language and virtual keyboard only \\
1 & Top Bar        & Drag-and-drop layout of the Top Bar instrument strip & No \\
2 & Bottom Bar     & Drag-and-drop layout of the Bottom Bar action strip & No \\
3 & Widgets        & Define and edit reusable data widget definitions & No \\
4 & Comfort        & Select and customise visual comfort presets & No \\
5 & Connection     & Signal~K server URL, authentication token, mDNS discovery & No \\
6 & Signal~K       & Map semantic names to actual Signal~K paths & No \\
7 & Units          & Choose display units per measurement category & No \\
8 & Applications   & Manage the launcher: pages, tiles, folders, app details & No \\
9 & System         & Diagnostics, logging, performance, configuration management & No (some actions restart) \\
10 & Android       & Optional Home role, native launcher application selection, soft navigation controls & No (Android only) \\
\bottomrule
\end{longtable}

On Android, the Android tab lists only launchable \code{MAIN + LAUNCHER}
actividades. Selección de una aplicación la añade a la paleta de aplicaciones compartidas sin
solicitando una amplia visibilidad del paquete. Identidad de aplicación nativa, icono, paquete y
actividad son gestionados por Android. Deseleccionar una aplicación elimina su página de lanzamiento
referencias. FairWindSK nunca se convierte en la aplicación Home predeterminada automáticamente, y la
La intención del lanzador ordinario sigue disponible. Controles de espalda suave, hogar y recientes
aparecen sólo cuando Android informa que las teclas de hardware correspondientes están ausentes;
Recents contains only native activities launched through FairWindSK.

% ── Chapter 21: Tab 0 Main ──────────────────────────────────
\chapter{Ajustes --- Main}
\label{ch:settings-main}
\settingstab{0}{Main}

La pestaña Principal controla el comportamiento básico del tiempo de ejecución: modo de visualización de ventanas, idioma, elemento UI
escala, preset de confort activo, dimensiones de la red de lanzamiento y formato de coordenadas.

\screenshot{settings-main}{Ajustes --- Ficha principal: modo de ventana, idioma, escala UI, preset de confort, rejilla de lanzamiento, formato de coordenadas}{}

\section{Modo de ventana}

\rowcolors{2}{LtGray}{white}
\begin{longtable}{L{2.4cm}L{5.6cm}L{5.7cm}}
\toprule
\rowcolor{Navy}
\color{white}\sffamily\bfseries Mode &
\color{white}\sffamily\bfseries Behaviour &
\color{white}\sffamily\bfseries Notes \\
\midrule
Ventana&Ventana de escritorio removible en la izquierda/Top/Width/Height configurada.&
  Development, testing, multi-monitor setups. \\
Centred&Ventana centrada de la anchura configurada y la altura. Campos de posición discapacitados.&
  Preferred windowed mode without an exact position. \\
Maximizado&Área de trabajo de escritorio completa.&
En Raspberry Pi~OS, utiliza la geometría del área de trabajo de pantalla directamente,
  bypassing the window manager. \\
Pantalla completa&Ventana de kiosco sin marco, siempre en la parte superior.&
Recomendado para pantallas de helm dedicadas.
  On labwc desktops, set \code{autohide=true}
  in \code{wf-panel-pi.ini} if the panel overlays FairWindSK. \\
\bottomrule
\end{longtable}

\section{Idioma}

FairWindSK barcos con cuatro idiomas de interfaz totalmente compatibles.
La configuración del lenguaje afecta a todo el texto de cáscara nativo --- etiquetas de barras, páginas de configuración,
drawers, warnings, and status messages --- as well as the \code{Accept-Language}
header sent to embedded web content.

\begin{itemize}
  \item \textbf{System default} --- follows the OS locale when FairWindSK supports
ese idioma; vuelve al inglés por cada lenguaje OS no compatible.
  \item \textbf{English} --- forces English regardless of OS locale.
  \item \textbf{Français} --- forces French.
  \item \textbf{Español} --- forces Spanish.
  \item \textbf{Italiano} --- forces Italian. All native shell text is displayed in Italian.
\end{itemize}

\subsection{Cómo funciona la localización}

The \code{Localization} module normalises the selected language, chooses the
effective \code{QLocale}, and builds the web \code{Accept-Language} header.
The Qt translator (\code{.qm} resource) is loaded during startup, so native widgets
y las vistas web incrustadas comienzan con el mismo idioma y cultura.
Translation source files live in \code{resources/i18n/} (for example,
\code{fairwindsk\_fr.ts}, \code{fairwindsk\_es.ts}, and
\code{fairwindsk\_it.ts}).

\subsection{Añadiendo o mejorando un idioma}

See \code{docs/multilingual.md} in the repository for the full contributor workflow:
adding translatable strings with \code{tr(\ldots)}, updating \code{.ts} files with
Qt Linguist, registering a new language in \code{Localization.cpp} and
\code{ui/settings/Main.cpp}, and updating \code{CMakeLists.txt}.

\begin{warningbox}
Se requiere un reinicio después de cambiar el idioma.
FairWindSK muestra un cajón informativo cuando se cambia el idioma.
Sin reiniciar, los widgets nativos Qt y la vista web incrustada pueden utilizar diferentes
idiomas.
\end{warningbox}

\section{UI Scale Preset}

Cuatro presets: Pequeño, Normal (por defecto), Grande, Extra Grande.
When \key{Auto Scale} is ticked, FairWindSK selects the preset based on display resolution
y DPI.

\section{Vista Comfort Preset y conmutación automática}

Selecciona el preset de confort activo directamente desde la pestaña Principal.
When \key{Auto Comfort View} is enabled, FairWindSK switches presets automatically
using Signal~K sun-position data (\skpath{environment.sun}).
Ambas condiciones deben cumplirse para que el modo auto esté disponible:

\begin{enumerate}
  \item The \skpath{environment.sun} path must be configured in the Signal~K tab.
  \item The Signal~K server must be connected and delivering values on that path.
\end{enumerate}

\section{Lanzador Grid}

Dos cajas de giro fijan el número de filas y columnas de fichas de aplicación por página de lanzamiento.
Menos filas y columnas producen azulejos más grandes que son más fáciles de tocar.
Más producir más aplicaciones visibles por página.

\section{Coordinate Format}

Tres formatos están disponibles a lo largo de la interfaz (propulsor de posición de la barra, barra de POB,
Anchor bar, My~Data waypoints:

\rowcolors{2}{LtGray}{white}
\begin{longtable}{L{3.2cm}L{4.8cm}L{5.7cm}}
\toprule
\rowcolor{Navy}
\color{white}\sffamily\bfseries Format &
\color{white}\sffamily\bfseries Example &
\color{white}\sffamily\bfseries When to use \\
\midrule
DD.DDDDD\textdegree{} & 48.85341\textdegree{} N & GPS logs, digital charting software. \\
DD\textdegree{}MM.MMM' & 48\textdegree{}51.205' N & Standard nautical format. Recommended for most sailors. \\
DD\textdegree{}MM'SS.S'' & 48\textdegree{}51'12.3'' N & Traditional celestial navigation format. \\
\bottomrule
\end{longtable}

\section{Teclado virtual}

Permite el teclado virtual Qt para implementaciones de pantalla táctil.
Se requiere un reinicio después de regatear.
En Raspberry~Pi~OS, el script de ayuda de arranque aplica la variable de entorno correcto
(\code{QT\_IM\_MODULE=qtvirtualkeyboard}) automatically.

% ── Chapter 22: Tabs 1 and 2 --- Bar Editors ──────────────────
\chapter{Ajustes --- Top Bar and Bottom Bar Editors}
\label{ch:settings-bars}
\settingstab{1}{Top Bar} \settingstab{2}{Bottom Bar}

Tabs~1 y~2 comparten el mismo editor visual de arrastrar y soltar --- uno para el Top~Bar y uno
Para el Bottom~Bar.

\screenshot{settings-topbar-editor}{Ajustes --- Top Bar editor: cinta de vista previa en vivo, barra de herramientas de edición y paleta widget}{}

\section{Vista previa en vivo}

La vista previa en la parte superior hace una simulación de ancho completo de la barra utilizando la comodidad actual
preajustar colores y datos reales de widget.
Toque cualquier artículo en la vista previa para seleccionarlo.

\section{Edición de botones de barra de herramientas}

These are the exact icons defined in \code{BarLayoutSettings.cpp}.

\rowcolors{2}{LtGray}{white}
\begin{longtable}{C{1.5cm}L{4.0cm}L{8.2cm}}
\toprule
\rowcolor{Navy}
\color{white}\sffamily\bfseries Icon &
\color{white}\sffamily\bfseries Button &
\color{white}\sffamily\bfseries Action \\
\midrule
\iconlg{layout-horizontal-minimize} & Minimize width  &
  Widget uses only the space its content requires. \\
\iconlg{layout-horizontal-maximize} & Maximize width  &
  Widget expands to fill all available space horizontally. \\
\iconlg{layout-vertical-minimize}   & Minimize height &
  Widget uses the standard bar height. \\
\iconlg{layout-vertical-maximize}   & Maximize height &
  Widget uses the full bar height (both header and value rows fully visible). \\
\iconlg{arrow-left-google}          & Move left       &
  Moves the selected widget one position to the left. \\
\iconlg{arrow-right-google}         & Move right      &
  Moves the selected widget one position to the right. \\
\iconlg{delete-google}              & Remove selected &
  Returns the selected widget to the palette. \\
\iconlg{refresh-google}             & Reset defaults  &
  Restores the bar to its built-in default layout. \\
\bottomrule
\end{longtable}

\section{Pantalla Toggles}

Con un widget seleccionado, cuatro casillas controlan lo que se muestra:
\key{Show Icon}, \key{Show Text} (label), \key{Show Units}, \key{Show Trend} (arrow).

\section{Widget Palette}

La paleta debajo de la barra de herramientas lista todos los elementos disponibles.
\textbf{Data widgets} are backed by Signal~K paths.
\textbf{Fixed shell controls} --- Apps, POB, Autopilot, Anchor, Alarms, My~Data, Settings,
Signal~K status, Clock --- se puede mover entre el Top~Bar y Bottom~Bar.
\textbf{Separators} (\iconlg{layout-separator}) and \textbf{Stretches}
(\iconlg{layout-elastic-extender}) are invisible layout helpers.

% ── Chapter 23: Tab 3 --- Widgets ─────────────────────────────
\chapter{Ajustes --- Widgets}
\label{ch:settings-widgets}
\settingstab{3}{Widgets}

La pestaña Widgets gestiona la biblioteca de definiciones widget de datos utilizadas por ambas barras.

\screenshot{settings-widgets}{Ajustes --- pestaña Widgets con widget lista y campos de definición}{}

\section{Campo de definición de Widget}

\rowcolors{2}{LtGray}{white}
\begin{longtable}{L{2.8cm}L{10.9cm}}
\toprule
\rowcolor{Navy}
\color{white}\sffamily\bfseries Field &
\color{white}\sffamily\bfseries Purpose \\
\midrule
id            & Unique identifier used internally and in bar layout entries. \\
name          & Human-readable label shown in bars and this editor. \\
icon          & Optional SVG icon path. Accepts Qt resource paths or filesystem paths. \\
type          & Rendering format: \code{numerical}, \code{gauge}, \code{position},
                \code{datetime}, or \code{waypoint}. \\
signalKPath   & Full dot-separated Signal~K data path. \\
sourceUnit    & Unit in which Signal~K delivers the value (e.g.\ \code{m/s}, \code{rad}). \\
defaultUnit   & Unit to display in (e.g.\ \code{kn}, \code{deg}).
                Overridden by active unit preferences. \\
updatePolicy  & \code{instant} (default) or \code{fixed}. \\
period        & Target update period in milliseconds (default 1000). \\
minPeriod     & Minimum update period in milliseconds (default 200). \\
\bottomrule
\end{longtable}

% ── Chapter 24: Tab 4 --- Comfort ─────────────────────────────
\chapter{Ajustes --- Comfort}
\label{ch:settings-comfort}
\settingstab{4}{Comfort}

La pestaña Comfort es el editor de temas visuales para los seis presets de confort incorporados.

\screenshot{settings-comfort}{Ajustes --- pestaña Comfort: selector de presets, relojes de color y sección de imagen de fondo}{}

\section{Los Six Comfort Presets}

\rowcolors{2}{LtGray}{white}
\begin{longtable}{C{1.5cm}L{2.0cm}L{4.2cm}L{6.0cm}}
\toprule
\rowcolor{Navy}
\color{white}\sffamily\bfseries Icon &
\color{white}\sffamily\bfseries Preset &
\color{white}\sffamily\bfseries Conditions &
\color{white}\sffamily\bfseries Key characteristics \\
\midrule
\iconlg{comfort-default} & Default & General fallback &
  Balanced mid-contrast. Usable in most conditions. \\
\iconlg{comfort-dawn}    & Dawn    & Around sunrise &
  Warm tones, slightly reduced brightness. Good contrast. \\
\iconlg{comfort-day}     & Day     & Full daylight, direct sun &
  Maximum brightness and contrast. For open-ocean sailing. \\
\iconlg{comfort-sunset}  & Sunset  & Around sunset &
  Warmer, transitional brightness. \\
\iconlg{comfort-dusk}    & Dusk    & Twilight &
  Reduced brightness, begins night-vision preservation. \\
\iconlg{comfort-night}   & Night   & Full dark, offshore watches &
Brillo mínimo, fondo oscuro, acentos rojos.
  Preserves night vision. Alarms remain high-contrast. \\
\bottomrule
\end{longtable}

\section{Botones de acción}

\rowcolors{2}{LtGray}{white}
\begin{longtable}{C{1.5cm}L{3.0cm}L{9.2cm}}
\toprule
\rowcolor{Navy}
\color{white}\sffamily\bfseries Icon &
\color{white}\sffamily\bfseries Button &
\color{white}\sffamily\bfseries Action \\
\midrule
\iconlg{anchor-reset}     & Reset       &
  Clears all custom colour overrides, background images, and QSS for the selected preset. \\
\iconlg{delete-google}    & Reset All   &
  Clears customisation for all six presets. \\
\iconlg{route-export-iec} & Export QSS  &
  Saves the selected preset's effective stylesheet to a \code{.qss} file. \\
\iconlg{route-import-iec} & Import QSS  &
  Loads a \code{.qss} file as the custom stylesheet for the selected preset. \\
\bottomrule
\end{longtable}

\section{Propiedades de anulación de color}

Diez propiedades de color son personalizables por preset:

\rowcolors{2}{LtGray}{white}
\begin{longtable}{L{3.6cm}L{10.1cm}}
\toprule
\rowcolor{Navy}
\color{white}\sffamily\bfseries Property &
\color{white}\sffamily\bfseries What it controls \\
\midrule
Window background  & Background colour of the Application Area and settings pages. \\
Text               & Primary text colour throughout the interface. \\
Button background  & Background of buttons and tool buttons in bars and drawers. \\
Button text        & Text and icon tint on buttons. \\
Border             & Borders, dividers, and frame outlines. \\
Accent             & Primary accent colour for highlights and selected states. \\
Accent text        & Text colour on accent-coloured backgrounds. \\
Icon               & Default tint for monochrome SVG icons. Set to red for Night preset to aid night vision. \\
Scroll track       & Background of scroll bar tracks. \\
Scroll knob        & Colour of the scroll bar handle. \\
\bottomrule
\end{longtable}

\begin{warningbox}
Nunca configure un preset que haga texto de alarma, información de POB o crítica de seguridad
lecturas difíciles de leer.
La lectura en todas las condiciones tiene precedencia sobre la estética.
\end{warningbox}

% ── Chapter 25: Tab 5 --- Connection ──────────────────────────
\chapter{Ajustes --- Conexión}
\label{ch:settings-connection}
\settingstab{5}{Connection}

\screenshot{settings-connection}{Ajustes --- Ficha de conexión: campo URL del servidor, línea de estado, botones de acción y registro de consola}{}

\section{Campo de URL del servidor}

An editable combo box pre-populated with the configured URL, \code{http://localhost:3000},
and \code{https://demo.signalk.org}.
FairWindSK normalises the URL: prepends \code{http://} if no scheme is present,
strips trailing slashes, and validates that only \code{http://} or \code{https://} are accepted.
Typed URLs are committed to the saved configuration only when you tap \key{Connect}.

\section{Status Line}

Una sola línea debajo del campo URL muestra el estado actual de token y conexión utilizando
three dot-separated fields: \textbf{State}, \textbf{Permission}, and \textbf{Expires}.

\section{Botones de acción}

\subsection{Conectar y Pausa}

Toggles the Signal~K connection.
\key{Connect} starts the connection; \key{Pause} suspends it without clearing the URL.

\subsection{Solicitud Token}

Inicia el flujo de autenticación de Signal~K:

\begin{enumerate}
  \item FairWindSK posts a new access request to \skpath{/signalk/v1/access/requests}.
  \item On HTTP~202 + \code{PENDING}, FairWindSK opens the Signal~K admin interface in the
Cajón de la barra inferior, navegado a la seguridad~ Página de peticiones de acceso.
  \item FairWindSK polls the \code{href} every second until resolved.
  \item On \code{COMPLETED} with \code{APPROVED}, the token is stored in
        \code{fairwindsk.ini} and injected as the \code{JAUTHENTICATION} cookie in the
tienda de galletas del navegador incrustado.
\end{enumerate}

\subsection{Cancelar}

Aborta la solicitud de señal de vuelo y aclara el href pendiente.

\subsection{Quitar Token}

Removes the stored token from memory, \code{fairwindsk.ini}, and the browser cookie store.
Provoca una reconexión en el modo público (no autenticado).

\section{Consola Log}

Un área de texto de sólo lectura desplazable debajo de los botones muestra eventos de conexión:
mDNS mensajes de descubrimiento, eventos del ciclo de vida de conexión y resultados de petición de fichas.

\section{mDNS Automatic Discovery}

En plataformas de escritorio, FairWindSK escanea la red local para servidores Signal~K usando
mDNS/Zeroconf en cuatro tipos de servicio:
\code{\_signalk-http.\_tcp}, \code{\_signalk-https.\_tcp}, \code{\_http.\_tcp},
and \code{\_https.\_tcp}.
Los servidores descubiertos se añaden automáticamente al cuadro combo URL.
mDNS está deshabilitado en Android e iOS construye.

% ── Chapter 26: Tab 6 --- Signal K Paths ──────────────────────
\chapter{Ajustes --- Signal K Paths}
\label{ch:settings-signalk}
\settingstab{6}{Signal K Paths}

Maps the semantic names that FairWindSK uses internally to the actual Signal~K data paths
en su servidor.
Editar cualquier campo para cambiar donde FairWindSK lee o escribe ese valor.
Los cambios tienen efecto después de un breve retraso en el desembolso sin necesidad de reiniciar.

\screenshot{settings-signalk}{Ajustes --- Señal K Paths pestaña con campos de etiqueta y ruta editable}{}

\section{Mappings de sendero clave}

\rowcolors{2}{LtGray}{white}
\begin{longtable}{L{3.0cm}L{5.6cm}L{5.1cm}}
\toprule
\rowcolor{Navy}
\color{white}\sffamily\bfseries Key &
\color{white}\sffamily\bfseries Default Signal~K path &
\color{white}\sffamily\bfseries Feature \\
\midrule
\code{pos}              & \skpath{navigation.position}                              & Top~Bar position, POB position \\
\code{sog}              & \skpath{navigation.speedOverGround}                       & Top~Bar SOG widget \\
\code{cog}              & \skpath{navigation.courseOverGroundTrue}                   & Top~Bar COG widget \\
\code{hdg}              & \skpath{navigation.headingTrue}                           & Top~Bar HDG widget \\
\code{dpt}              & \skpath{environment.depth.belowTransducer}                & Top~Bar depth widget \\
\code{notifications.pob}& \skpath{notifications.mob}                                & POB alarm \\
\code{anchor.position}  & \skpath{navigation.anchor.position}                      & Anchor Watch \\
\code{anchor.radius}    & \skpath{navigation.anchor.currentRadius}                 & Anchor Watch \\
\code{anchor.max}       & \skpath{navigation.anchor.maxRadius}                     & Anchor Watch \\
\code{anchor.rode}      & \skpath{winches.windlass.rode}                           & Anchor rode display \\
\code{anchor.actions.drop}   & \skpath{plugins.anchoralarm.dropAnchor}          & Anchor Drop button \\
\code{anchor.actions.raise}  & \skpath{plugins.anchoralarm.raiseAnchor}         & Anchor Raise button \\
\code{anchor.actions.radius} & \skpath{plugins.anchoralarm.setRadius}           & Anchor Set Radius button \\
\code{anchor.actions.up}     & \skpath{plugins.windlassctl.up}                  & Anchor Up button \\
\code{anchor.actions.down}   & \skpath{plugins.windlassctl.down}                & Anchor Down button \\
\code{environment.sun}  & \skpath{environment.sun}                                  & Auto comfort preset switching \\
\code{autopilot.state}  & \skpath{steering.autopilot.state}                        & Autopilot bar mode display \\
\code{autopilot.mode}   & \skpath{steering.autopilot.mode}                         & Autopilot engage commands \\
\code{rsa}              & \skpath{steering.rudderAngle}                            & Autopilot rudder indicator \\
\bottomrule
\end{longtable}

% ── Chapter 27: Tab 7 --- Units ────────────────────────────────
\chapter{Ajustes --- Unidades}
\label{ch:settings-units}
\settingstab{7}{Units}

\screenshot{settings-units}{Ajustes --- Ficha de unidades: nombre predefinido del servidor y filas de anulación por categoría}{}

\section{Cómo funcionan las preferencias de la unidad}

FairWindSK lee las preferencias de la unidad del servidor Signal~K conectado.
Cada fila de categoría muestra: la etiqueta de categoría, una caja combo de anulación (su preferencia local),
y el valor actual del servidor (sólo lectura).
Si elige la misma unidad que el servidor, no se almacena la anulación.
If different, it is saved in \code{unitPreferences.overrides.CATEGORY} in
\code{fairwindsk.json} and takes precedence.

% ── Chapter 28: Tab 8 --- Applications ────────────────────────
\chapter{Ajustes --- Aplicaciones}
\label{ch:settings-apps}
\settingstab{8}{Applications}

Un editor de dos pane: el panel izquierdo administra las páginas de lanzamiento y carpetas; el panel derecho es
la paleta de aplicaciones.

\screenshot{settings-apps}{Ajustes --- Ficha de aplicaciones: árbol de página (izquierda) y paleta de aplicaciones (derecha)}{}

\section{Acciones del árbol de página}

\rowcolors{2}{LtGray}{white}
\begin{longtable}{C{1.5cm}L{3.8cm}L{8.4cm}}
\toprule
\rowcolor{Navy}
\color{white}\sffamily\bfseries Icon &
\color{white}\sffamily\bfseries Action &
\color{white}\sffamily\bfseries What it does \\
\midrule
\iconlg{add-new-page}                    & Add page       & Creates a new empty launcher page. \\
\iconlg{show-page-details}               & Page details   & Opens the page details drawer (name, icon). \\
\iconlg{add-all-apps-to-current-page}    & Assign all apps & Adds all palette apps to the current page. \\
\iconlg{remove-all-apps-from-current-page}& Clear page    & Removes all tiles from the current page. \\
\iconlg{remove-all-apps-from-all-pages}  & Clear all pages & Removes all tiles from every page. \\
\bottomrule
\end{longtable}

\section{Palette Actions}

\rowcolors{2}{LtGray}{white}
\begin{longtable}{C{1.5cm}L{4.0cm}L{8.2cm}}
\toprule
\rowcolor{Navy}
\color{white}\sffamily\bfseries Icon &
\color{white}\sffamily\bfseries Action &
\color{white}\sffamily\bfseries What it does \\
\midrule
\iconlg{define-new-application-in-palette}  & Create app    & Opens the App Details drawer to create a new custom application. \\
\iconlg{show-details-selected-app-in-palette} & App details & Opens the App Details drawer to edit the selected application. \\
\iconlg{remove-selected-app-from-palette}   & Remove app    & Permanently removes the selected application from the palette. \\
\iconlg{add-selected-app-to-current-page}   & Add to page   & Adds the selected palette app to the current launcher page. \\
\bottomrule
\end{longtable}

\section{Detalles de la aplicación}

\rowcolors{2}{LtGray}{white}
\begin{longtable}{L{2.8cm}L{10.9cm}}
\toprule
\rowcolor{Navy}
\color{white}\sffamily\bfseries Field &
\color{white}\sffamily\bfseries Description \\
\midrule
Display name    & Name shown on the launcher tile and in the Top Bar when the app is active. \\
URL             & Full URL (\code{http://} or \code{https://}). \\
Icon path       & Qt resource path, filesystem path, or bundled icon path. \\
Web zoom level  & Percentage (default 100\%). Increase for larger buttons; decrease for more content. \\
Active          & Whether the tile is visible on the launcher. \\
Order           & Secondary sort hint within a page. \\
\bottomrule
\end{longtable}

% ── Chapter 29: Tab 9 --- System ──────────────────────────────
\chapter{Ajustes --- Sistema}
\label{ch:settings-system}
\settingstab{9}{System}

The System tab is divided into two panes: \textbf{Performance} (left) and
\textbf{Diagnostics and Logging} (right).

\screenshot{settings-system-perf}{Ajustes --- Ficha del sistema Pane de rendimiento: barras de memoria, CPU por núcleo, red, métricas Raspberry Pi}{}

\section{Performance Pane}

\subsection{Memory Readouts}

\begin{itemize}
  \item \textbf{Process memory} --- resident RAM used by FairWindSK
        (from \code{/proc/self/statm} on Linux; Mach task API on macOS).
  \item \textbf{Total memory} --- installed RAM
        (from \code{/proc/meminfo} on Linux; \code{sysctl hw.memsize} on macOS).
  \item \textbf{Memory usage bar} --- process memory as a percentage of total RAM.
\end{itemize}

\subsection{Carga CPU}

Una fila de barras horizontales de progreso, una por núcleo de CPU, actualizada cada segundo.
On Linux, data is read from \code{/proc/stat};
on macOS from the Mach \code{host\_processor\_info} API.

\subsection{Direcciones de red}

Listas direcciones IPv4 de todas las interfaces de red activas y no retrocesivas, refrescada cada 10 segundos.
Útil para confirmar qué dirección IP deben usar otros dispositivos para llegar a esta instalación.

\subsection{Diagnósticos de fresa Pi}

Cuando el servidor Signal~K conectado expone las métricas Raspberry~Pi
(via the \code{signalk-server-raspberrypi} plugin on path \skpath{environment.rpi}),
un grupo Raspberry~Pi aparece mostrando:

\rowcolors{2}{LtGray}{white}
\begin{longtable}{L{3.6cm}L{10.1cm}}
\toprule
\rowcolor{Navy}
\color{white}\sffamily\bfseries Metric &
\color{white}\sffamily\bfseries Description \\
\midrule
CPU temperature   & SoC CPU temperature in \textdegree{}C. Above 80\textdegree{}C indicates thermal throttling. \\
GPU temperature   & VideoCore GPU temperature in \textdegree{}C. \\
CPU utilisation   & Overall CPU load as a percentage. \\
Memory utilisation& RAM usage as a percentage. \\
SD utilisation    & Storage usage on the SD card as a percentage. \\
\bottomrule
\end{longtable}

\section{Diagnósticos y Panes de Logging}

\screenshot{settings-system-log}{Ajustes --- Ficha del sistema Diagnósticos pane: nivel de registro, registros persistentes, historia de interacción, caminos, ver registros y ver los informes botones}{}

\subsection{Nivel de registro}

\rowcolors{2}{LtGray}{white}
\begin{longtable}{L{2.4cm}L{4.6cm}L{6.7cm}}
\toprule
\rowcolor{Navy}
\color{white}\sffamily\bfseries Level &
\color{white}\sffamily\bfseries Messages included &
\color{white}\sffamily\bfseries When to use \\
\midrule
Off (default) & No messages. & Normal operation. Minimises SD card writes on Raspberry~Pi. \\
Critical      & Fatal errors only. & Minimal logging with crash detection. \\
Warning       & Critical + non-fatal anomalies. & Initial installation testing. \\
Info          & Warning + lifecycle events. & Commissioning and setup. \\
Debug         & Info + detailed internal state. & Troubleshooting. Significant overhead. \\
Full          & All of the above + verbose tracing. & Developer use only. Very high overhead. \\
\bottomrule
\end{longtable}

\subsection{Registros persistentes}

Cuando está habilitado, FairWindSK escribe un registro completo de mensajes para cada carrera al usuario
directorio de datos de aplicación, con el timetamp de inicio de proceso en el nombre de archivo.

\subsection{Historia de la interacción}

Cuando está habilitado, FairWindSK escribe una revista ligera de navegación y control del operador
acciones junto a cada registro de ejecución.
Útil para la revisión posterior al incidente.

\subsection{Ver registros y ver informes Botones}

Abra el explorador de registro de aplicación en el cajón horizontal Bottom Bar.
Los archivos de registro aparecen a la izquierda; las entradas de registro pares aparecen a la derecha.
Crash-report packs (creados automáticamente cuando FairWindSK no salió limpiamente) son
accessible via \key{View Reports}.

\section{Botones de gestión de configuración}

\rowcolors{2}{LtGray}{white}
\begin{longtable}{L{3.2cm}L{10.5cm}}
\toprule
\rowcolor{Navy}
\color{white}\sffamily\bfseries Button &
\color{white}\sffamily\bfseries Action \\
\midrule
Reset&
Discuta cambios sin salvar y recarga la configuración activa.
  FairWindSK keeps running. \\
Restore Defaults&
Reemplaza todas las configuraciones con los predeterminados agrupados.
  All customisation is lost. Confirmation required. \\
Importación&
  Imports a \code{fairwindsk.json} file from the filesystem.
  Validates, atomically replaces the active configuration, and reloads immediately. \\
Exportación&
  Copies the current \code{fairwindsk.json} to a chosen location. \\
\bottomrule
\end{longtable}

\begin{warningbox}
\key{Restore Defaults} is irreversible without a backup.
Exporte su configuración primero si desea restaurarla.
\end{warningbox}

\section{Botones de control de sistemas}

\rowcolors{2}{LtGray}{white}
\begin{longtable}{L{3.2cm}L{10.5cm}}
\toprule
\rowcolor{Navy}
\color{white}\sffamily\bfseries Button &
\color{white}\sffamily\bfseries Action \\
\midrule
Restart FairWindSK&
Ahorra configuración, luego cierra FairWindSK con código de salida~1.
  The Raspberry~Pi~OS startup helper relaunches on exit code~1. \\
Quit FairWindsK&
Guarda la configuración y cierra limpiamente (código de salida~0).
  The startup helper does \textbf{not} relaunch after a clean exit. \\
\bottomrule
\end{longtable}

\begin{warningbox}
Quitting FairWindSK en una pantalla de helm dedicada deja la pantalla en blanco.
Asegurar que otro miembro de la tripulación mantenga un reloj seguro antes de dejar de fumar.
\end{warningbox}
